\documentclass{article}
\usepackage{amsmath}
\usepackage{siunitx}

\begin{document}

\section*{Ferrite Core Selection Using the \( K_g \) Method}

\subsection*{Input Parameters}
\begin{align*}
L &= \SI{34.9e-6}{\henry} & &\text{(Inductance)} \\
I_o &= \SI{4.583}{\ampere} & &\text{(DC Current)} \\
\Delta I &= \SI{1.375}{\ampere} & &\text{(AC Ripple Current)} \\
P_o &= \SI{110}{\watt} & &\text{(Output Power)} \\
\alpha &= 1\% & &\text{(Regulation)} \\
f_s &= \SI{100e3}{\hertz} & &\text{(Switching Frequency)} \\
B_m &= \SI{0.25}{\tesla} & &\text{(Operating Flux Density)} \\
K_u &= 0.4 & &\text{(Window Utilization)} \\
T_r &= \SI{25}{\celsius} & &\text{(Temperature Rise Goal)}
\end{align*}

\subsection*{Calculations}
\begin{enumerate}
\item \textbf{Peak Current (\( I_{\text{pk}} \))}:
\[
I_{\text{pk}} = I_o + \frac{\Delta I}{2} = 4.583 + \frac{1.375}{2} = \SI{5.27}{\ampere}
\]

\item \textbf{Energy-Handling Capability (\( E \))}:
\[
E = \frac{1}{2} L I_{\text{pk}}^2 = \frac{1}{2} \cdot 34.9 \times 10^{-6} \cdot (5.27)^2 = \SI{4.87e-4}{\joule}
\]

\item \textbf{Electrical Condition Coefficient (\( K_e \))}:
\[
K_e = 0.145 \cdot P_o \cdot B_m^2 \cdot 10^{-4} = 0.145 \cdot 110 \cdot (0.25)^2 \cdot 10^{-4} = \num{9.96e-5}
\]

\item \textbf{Core Geometry Coefficient (\( K_g \))}:
\[
K_g = \frac{E^2}{K_e \cdot \alpha} = \frac{(4.87 \times 10^{-4})^2}{9.96 \times 10^{-5} \cdot 1} = \num{2.38e-3}
\]
\end{enumerate}

\subsection*{Notes}
\begin{itemize}
\item \( K_g \) determines the minimum core size required
\item Select a ferrite core with \( K_g \geq \num{2.38e-3} \, \text{cm}^5 \)
\item Verify core saturation at \( B_m = \SI{0.25}{\tesla} \)
\end{itemize}

\end{document}